\documentclass{article}
\usepackage{graphicx}
\usepackage{amsmath}

\title{The Evolution of Computer Vision: From Pixels to Perception}
\author{Samarpan Verma}
\date{2026-06-12}

\begin{document}
\maketitle

\section{Introduction}

Computer Vision (CV) has undergone a dramatic transformation over the past decade. What once relied heavily on hand-crafted feature extractors like SIFT and HOG has now been entirely subsumed by deep learning architectures. This editorial explores the paradigm shifts in visual recognition systems and experiments with how we present image-heavy content in a single-column format.

The primary goal of a computer vision system is to map an image $I \in \mathbb{R}^{H \times W \times 3}$ to a specific semantic understanding, whether that is a class label $y \in \mathcal{Y}$ or a dense pixel-wise segmentation map.

\section{The Deep Learning Revolution}

Before 2012, image classification was constrained by the limits of human intuition in designing edge detectors and corner algorithms. The breakthrough came with Convolutional Neural Networks (CNNs), which autonomously learn hierarchical feature representations. 

\begin{figure}[h]
\centering
\includegraphics{https://images.unsplash.com/photo-1555949963-aa79dcee981c?q=80&w=800&auto=format&fit=crop}
\caption{The complexity of computer vision ranges from simple classification to abstract perception in complex scenes.}
\end{figure}

The core mathematical operation driving this revolution is the discrete convolution, defined for a 2D image $I$ and a kernel $K$ as:

\begin{equation}
(I * K)(i, j) = \sum_{m} \sum_{n} I(i - m, j - n) K(m, n)
\end{equation}

By stacking these convolutional layers, networks like AlexNet, VGG, and ResNet were able to capture increasingly complex concepts. ResNet introduced skip connections, which solved the vanishing gradient problem in extremely deep networks:

\begin{equation}
\mathcal{H}(x) = \mathcal{F}(x) + x
\end{equation}

where $\mathcal{H}(x)$ is the desired underlying mapping, and $\mathcal{F}(x)$ is the residual mapping to be learned.

\section{The Shift to Transformers}

While CNNs rely on local receptive fields, the Vision Transformer (ViT) treats an image as a sequence of patches, leveraging the self-attention mechanism originally designed for natural language processing. 

\begin{figure}[h]
\centering
\includegraphics{https://images.unsplash.com/photo-1620712943543-bcc4688e7485?q=80&w=800&auto=format&fit=crop}
\caption{Artificial neural networks have evolved to process vast amounts of unstructured visual data.}
\end{figure}

The self-attention mechanism computes the response at a position in a sequence by attending to all positions within the same sequence:

\begin{equation}
\text{Attention}(Q, K, V) = \text{softmax}\left(\frac{QK^T}{\sqrt{d_k}}\right)V
\end{equation}

This allows the network to build global context immediately, at the cost of a quadratic computational complexity with respect to the number of patches.

\section{Modern Applications}

Today, computer vision powers everything from autonomous driving to medical image analysis. Object detection systems like YOLO (You Only Look Once) approach the problem as a single regression task, directly predicting bounding boxes and class probabilities from full images in one evaluation.

\begin{figure}[h]
\centering
\includegraphics{https://images.unsplash.com/photo-1506443432602-ac2fcd6f54e0?q=80&w=800&auto=format&fit=crop}
\caption{Autonomous systems rely heavily on real-time computer vision inference.}
\end{figure}

\section{Conclusion}

As computational power continues to scale and architectures become increasingly multimodal, the boundary between textual understanding and visual perception is blurring. The future of computer vision lies in models that can not only identify objects but reason about their physical properties and relationships in a 3D world.

\end{document}
